% Compiles on its own (pulls ../main.tex's preamble via subfiles) or as part of the book.
\documentclass[../main.tex]{subfiles}
\begin{document}

\chapter{Training: The Loop}

Chapter 6 gave us a loss that turns and gradients that flow. This chapter is the theory gate for
the machine that runs that loss thousands of times: the training loop. It has two halves --- the
\emph{plumbing} that feeds the model batches (the training-ready dataset), and the \emph{control}
that drives learning (optimiser, learning rate, schedule, mixed precision, seeds). Neither is
exotic. But from-scratch training --- our regime, chosen because the model is small (8.94M) and the
data is small --- is fragile in exactly these knobs, so each earns its explanation.

\section{From loader to batch: the training-ready dataset}

The paired loader (Ch.~1) yields one frame at a time, at native size, with positives and ignore
regions mixed together, and both modalities present. The training step wants something narrower: a
\emph{stack} of equal-size images and a target list shaped exactly as the loss expects. Three jobs
bridge the gap; KAIST removes the hardest.

\begin{itemize}
  \item \textbf{Resize / letterbox --- not needed here.} A batch tensor cannot be ragged, so every
        image in a batch must share a size. Most datasets need resizing to enforce that. Every KAIST
        frame is already $640\times512$, and $640$ and $512$ are both divisible by $32$ --- exactly
        the total stride of the backbone's three downsampling stages. So every frame is already a
        legal, identical-size input. (When we later \emph{inject} scale misalignment, resizing
        returns. Not now.)
  \item \textbf{Target formatting.} The loss wants rows $[\text{cls}, x_1, y_1, x_2, y_2]$. The
        loader already stores boxes as $xyxy$ pixels. Two conversions remain: keep only the scored
        positives (the real \texttt{person} boxes), and set their class index to $0$, because the
        model has one class and it lives at index $0$.
  \item \textbf{Collate.} PyTorch's default batcher tries to stack \emph{everything} and chokes on
        the boxes, because one frame has two pedestrians and the next has five. So we write a tiny
        collate function: stack the images (all legal, all $640\times512$) into one
        $(B,C,H,W)$ tensor, and leave the targets as a plain length-$B$ list. That list-of-tensors
        is precisely what the loss loops over, one image at a time.
\end{itemize}

The same layer picks a single modality (this is the RGB-only / thermal-only baseline) and divides
pixels by $255$ to put them in $[0,1]$ for stable from-scratch training.

\begin{lesson}
For this baseline we \emph{drop} ignore regions from the targets entirely. That is a deliberate
simplification, and it is not perfectly correct: a grid cell sitting inside an ignore box (a
cyclist, a crowd blob, a \texttt{person?}) will be pushed toward ``background'' by the objectness
loss, even though we genuinely do not know what is there. The fully-correct handling masks the
objectness loss inside ignore regions so those cells are neither positive nor penalised --- the
same courtesy the \emph{evaluator} already extends them (Ch.~3). We add that mask before the
headline runs; it changes label quality, not the plumbing, so it is safe to defer past the
first loss curve.
\end{lesson}

\section{The optimiser, and one from-scratch subtlety}

We use \textbf{SGD with momentum} --- YOLOX's own choice --- not Adam. The update is ``step
downhill, with inertia so you coast through small bumps.'' Momentum $0.9$ and Nesterov are the
standard settings.

The one subtlety that matters more from scratch than from a checkpoint: we \textbf{exempt
BatchNorm parameters and biases from weight decay}. Weight decay gently pulls weights toward zero to
discourage overfitting --- sensible for the large convolution kernels, harmful for BN's scale/shift
and for biases, which must move freely to calibrate a network that began as pure noise. From a
pretrained model those are already calibrated and the choice barely registers; from scratch,
decaying them can stall early training. So we split parameters into two groups --- kernels get
decay, everything one-dimensional does not --- and hand the optimiser both.

\section{Learning rate and schedule: the shape that decides everything}

The learning rate is the step size, and it is the single most consequential number in the run. Too
large and the loss diverges --- and from scratch there is no good initialisation to catch the fall.
Too small and it crawls, wasting GPU-hours. Crucially, \textbf{there is no correct value known in
advance}: it depends on the model, the batch size, the data. From a pretrained model you could
borrow a known-good rate; from scratch you cannot. \emph{This is the whole reason the study sweeps
the learning rate before locking the recipe} --- it is an empirical unknown, so we measure it.

The rate does not stay fixed. It follows a two-phase shape:

\begin{itemize}
  \item \textbf{Warmup} (first $\sim$5\% of steps): ramp linearly from near zero up to the target
        rate. A freshly-initialised network emits garbage gradients in its first steps; a full-size
        rate there can knock it into a bad region it never leaves. Warmup lets it find its footing.
  \item \textbf{Cosine decay} (the rest): lower the rate smoothly to near zero along a cosine curve.
        Large steps early to explore the loss landscape, tiny steps late to settle into a minimum.
\end{itemize}

We implement this as a plain multiplier on the base rate, recomputed every step. The base rate is
the swept knob; the \emph{shape} stays fixed across all experiments.

\section{Mixed precision: why 12\,GB is enough}

Automatic Mixed Precision (AMP) runs most of the forward pass in 16-bit floats instead of 32-bit:
roughly half the memory and, on the laptop's Blackwell tensor cores, close to double the
throughput. That is the difference between a usable batch size and an out-of-memory crash. The one
hazard is that 16-bit floats underflow --- very small gradients round to zero. A \textbf{gradient
scaler} sidesteps this: multiply the loss up before the backward pass so gradients land in the
representable range, then divide it back out before the optimiser steps. It is boilerplate, but it
is the reason this model trains on this card at all.

\section{Seeds: the machinery behind ``5 seeds every experiment''}

One call at the top of the run seeds Python, NumPy, and Torch together. Same seed, same run. This
is not a nicety --- it is the foundation of the study's statistics. We run seeds $0$ through $4$ on
every experiment, report mean $\pm$ standard deviation, and treat a two-sigma band as the bar a
result must clear. Without reproducible seeds you cannot separate a real architectural improvement
from lucky initialisation, and the published KAIST literature --- which reports single-run deltas
of a few tenths of an MR point with no error bars anywhere --- is precisely the failure mode we
refuse to repeat.

\section{The exit check: watch it descend before you spend}

The chapter's discipline, and its exit check, is an ordering rule:

\begin{quote}
Before launching the $5$-seed $\times$ learning-rate sweep, run the loop once on a small subset and
\emph{watch the loss curve descend.}
\end{quote}

This is the ``overfit one image'' test of Ch.~6, scaled up to ``overfit a few hundred frames.'' If
the loss falls steadily on a subset of real KAIST data --- objectness dropping, IoU loss dropping,
the positive count staying sane --- then the loop, the dataset layer, the AMP path, and the
schedule are all wired correctly, and only then do we spend real GPU-hours on the sweep. Never burn
a $5$-seed budget on a loop you have not first watched turn.

\section{The first run failed --- and taught the chapter's real lesson}

The first smoke run did \emph{not} descend. Two things went wrong, and the second is the most
instructive failure in the project so far.

\textbf{Symptom one: \texttt{fg=0} on every step.} The subset was the first $300$ frames in sorted
order, which fell inside a stretch of pedestrian-free frames at the start of \texttt{set00/V000}.
SimOTA had nothing to assign because there were no pedestrians. This is why the smoke now filters
to frames that contain at least one scored person before taking its subset --- a plumbing test needs
something to learn.

\textbf{Symptom two: \texttt{loss=nan} from the moment warmup reached full learning rate.} Probing
one batch revealed the cause: the objectness loss at initialisation was \emph{eighteen thousand}.
At birth the head's output biases were zero, so every cell emitted a logit near $0$ and
$\sigma(0)=0.5$ --- the newborn model announced ``50\% chance of a pedestrian'' in all $6720$ cells.
The background penalty for that is $\approx 0.69$ per cell $\times\,6720\times B$ images $\approx
18{,}000$. That loss produced an enormous gradient; SGD at the target rate, amplified by the fp16
path, threw the weights to infinity within a few steps.

The fix is the \textbf{prior-probability bias initialisation} from RetinaNet, used by YOLOX itself.
Initialise the objectness and class output biases so that at birth every cell predicts ``almost
certainly background'':
\[
  b = -\log\frac{1-\pi}{\pi},\qquad \pi = 0.01 \;\Rightarrow\; b \approx -4.595,\quad \sigma(b)=0.01.
\]
Now each cell starts near the true object rate (a handful of pedestrians among thousands of cells),
the initial background loss falls roughly $50\times$, gradients stay tame, and warmup does its job.

\begin{lesson}
The overfit-one-image test of Ch.~6 was green, yet real training diverged. The test used
\emph{AdamW}, which divides each gradient by its own recent magnitude and so shrugged off an
$18{,}000$ loss and still learned. Real training uses \emph{SGD}, which applies the raw gradient,
and batches of mostly-background --- exactly the regime that exposes a missing prior-bias init. A
green suite proves the cases you tested, not correctness. The all-background, SGD, from-scratch case
was not in the suite; now it is, as a regression test asserting the initial background loss stays
small. When a unit test and a real run disagree, the real run is right --- find what the test's
easier conditions were hiding.
\end{lesson}

\section{The fixed loop descends}

With the prior-bias init in place (and a second, smaller fix: SimOTA's internal cross-entropy cost
must run with autocast disabled, because binary cross-entropy on already-sigmoided probabilities is
unsafe under mixed precision), the smoke was re-run on $300$ pedestrian-bearing frames for $20$
epochs. It descended cleanly:

\begin{center}
\begin{tabular}{lccccc}
\toprule
 & \texttt{loss} & \texttt{obj} & \texttt{iou} & \texttt{cls} & \texttt{fg} \\
\midrule
start (step 20)  & 42.2 & 37.4 & 0.70 & 1.35 & 13 \\
end (step 740)   & 1.10 & 0.16 & 0.12 & 0.36 & $\sim$85 \\
\bottomrule
\end{tabular}
\end{center}

Every reading is a health sign. The objectness loss started in the low tens --- not the eighteen
thousand of the broken run --- confirming the prior-bias fix, then fell $\sim$230$\times$ as the
model learned to silence the background. The box loss fell from $\approx 0.70$ (boxes barely
touching the pedestrian) to $\approx 0.12$ ($\approx$88\% IoU). The positive count \emph{rose} from
$13$ to $\sim$85: as predictions improve, more cells clear SimOTA's IoU bar and dynamic-$k$ hands
out more positives --- assignment and learning reinforcing each other. No NaN, and the learning-rate
schedule ramped to $0.01$ and cosine-decayed to zero exactly as designed.

\begin{lesson}
This is memorisation, not detection. Three hundred frames for twenty epochs with no validation set:
the model overfit a tiny slice, which is precisely what this gate is meant to prove --- that the
loop \emph{can} learn --- and nothing more. The low final loss says the machine works; it says
nothing about whether the detector generalises. That question belongs to the next phase: train on
the proper KAIST train split and report MR$^{-2}$ on the held-out test set.
\end{lesson}

\end{document}
